\documentclass[12pt]{article}
\usepackage[brazil]{babel}
\usepackage[utf8]{inputenc}
\usepackage[T1]{fontenc}
\usepackage{amsmath, amssymb}
\usepackage{geometry}
\usepackage{graphicx}
\usepackage{tikz}
\usetikzlibrary{shapes,arrows,positioning,calc,matrix,arrows.meta,decorations.pathmorphing}
\usepackage{pgfplots}
\pgfplotsset{compat=1.18}
\usepackage{xcolor}
\usepackage{enumitem}
\usepackage{caption}
\usepackage{float}

\definecolor{maincolor}{RGB}{0,100,50}
\usepackage{fancyhdr}
\definecolor{forestgreen}{RGB}{0,100,50}
\usepackage[most]{tcolorbox}

\geometry{margin=2.5cm}
\pagestyle{fancy}
\fancyhf{}

\renewcommand{\headrulewidth}{2pt}
\renewcommand{\footrulewidth}{2pt}

\renewcommand{\headrule}{\hbox to\headwidth{%
  \color{black}\leaders\hrule height \headrulewidth\hfill}}

\renewcommand{\footrule}{\hbox to\headwidth{%
  \color{black}\leaders\hrule height \footrulewidth\hfill}}

\setlength{\headheight}{51pt} 

\fancyhead[L]{%
  \includegraphics[height=1.2cm]{Universidade_de_Rio_Verde.jpg}
}

\fancyhead[R]{%
  \raggedleft
  \textcolor{black}{%
    \small
    \textbf{Lista de Exercícios --- Grafos e Árvores}\\
    Me.~Sergio Souza Novak\\
    \textit{Lógica para Engenharia de Software}
  }
}

\fancyfoot[C]{%
  \textcolor{black}{\thepage}
}

\newcommand{\quadro}[1]{\vspace{0.2cm}\begin{tcolorbox}[colback=white, colframe=black, height=#1, boxrule=0.5pt, arc=0pt] \textit{Resposta:} \end{tcolorbox}}

\begin{document}

\begin{enumerate}[label=\textbf{\arabic*.}]

% -------------------------------------------------------
% Questão 1: Grafos Completos (Kn)
% -------------------------------------------------------
\item \textbf{Grafos Completos ($K_n$): Definição e Propriedades}
Um grafo completo $K_n$ é aquele em que existe uma aresta conectando cada par de vértices distintos.
\begin{enumerate}[label=\alph*)]
    \item Desenhe o grafo $K_7$ no espaço abaixo.
    \item Qual é o grau de cada vértice em um $K_7$?
    \item Calcule o número total de arestas de um grafo completo $K_{10}$. Mostre a fórmula utilizada.
\end{enumerate}
\quadro{10cm}

% -------------------------------------------------------
% Questão 2: Ciclos Hamiltonianos
% -------------------------------------------------------
\item \textbf{Ciclos Hamiltonianos em Grafos Completos}
Sabe-se que um Ciclo Hamiltoniano é um caminho fechado que visita todos os vértices exatamente uma vez, voltando à origem.
\begin{center}
    \begin{tikzpicture}[scale=1.7, every node/.style={circle, draw=black, thick, minimum size=0.5cm, inner sep=0pt}]
        \node (1) at (90:1.2) {};
        \node (2) at (18:1.2) {};
        \node (3) at (306:1.2) {};
        \node (4) at (234:1.2) {};
        \node (5) at (162:1.2) {};
        % All edges (K5)
        \draw[gray!50] (1)--(3) (1)--(4) (2)--(4) (2)--(5) (3)--(5);
        % Hamilton cycle
        \draw[black, thick, ->] (1)--(2);
        \draw[black, thick, ->] (2)--(3);
        \draw[black, thick, ->] (3)--(4);
        \draw[black, thick, ->] (4)--(5);
        \draw[black, thick, ->] (5)--(1);
    \end{tikzpicture}
\end{center}
Utilizando a fórmula correta, calcule quantos ciclos Hamiltonianos distintos existem em um grafo completo de 6 vértices ($K_6$). Apresente a fórmula e o cálculo.
\quadro{4cm}

% -------------------------------------------------------
% Questão 5: Bipartido Completo
% -------------------------------------------------------
\item \textbf{Grafos Bipartidos Completos ($K_{m,n}$)}
Desenhe o grafo bipartido completo $K_{3,2}$, separando os conjuntos em esquerda e direita.
\quadro{5cm}

% -------------------------------------------------------
% Questão 6: Aplicação Prática Bipartidos
% -------------------------------------------------------
\item \textbf{Aplicações: O Problema da Atribuição}
Uma \textit{software house} possui 4 programadores (Ana, Bob, Carlos, Dani) e 4 tarefas urgentes (Frontend, Backend, IA, DevOps). Sabe-se que: Ana domina Frontend e IA; Bob sabe apenas Backend; Carlos sabe DevOps e IA; Dani sabe Frontend e DevOps.
\begin{enumerate}[label=\alph*)]
    \item Modele essa situação como um grafo bipartido, desenhando os nós e as arestas.
    \item Mostre, destacando no grafo ou listando os pares, se existe um \textbf{Casamento Perfeito} (\textit{Perfect Matching}), ou seja, um cenário onde cada programador receba exatamente 1 tarefa que sabe fazer sem que falte tarefa para ninguém.
\end{enumerate}
\quadro{6cm}

% -------------------------------------------------------
% Questão 7: Árvores - Classificação
% -------------------------------------------------------
\item \textbf{Árvores Binárias: Classificação}
Na ciência da computação, a estrutura das árvores importa muito para a performance dos algoritmos.
\begin{enumerate}[label=\alph*)]
    \item Qual a diferença fundamental entre uma Árvore Binária \textbf{Estrita} (Full) e uma Árvore Binária \textbf{Perfeita}?
    \item É possível que uma árvore seja considerada ``Completa'' mas não seja ``Perfeita''? Dê um exemplo desenhando uma pequena árvore de 4 nós que ilustre isso.
\end{enumerate}
\quadro{5cm}


% -------------------------------------------------------
% Questão 9: Teoremas de Altura
% -------------------------------------------------------
\item \textbf{Teoremas de Capacidade em Árvores}
Em Engenharia de Software, o estudo da complexidade de banco de dados frequentemente utiliza as propriedades expansivas de árvores binárias.
\begin{center}
    \begin{tikzpicture}[level distance=0.8cm,
        level 1/.style={sibling distance=3cm},
        level 2/.style={sibling distance=1.5cm},
        every node/.style={circle, draw=black, thick, minimum size=0.4cm, inner sep=0pt}]
        \node (root) {}
          child {node {}
            child {node {}}
            child {node {}}
          }
          child {node {}
            child {node {}}
            child {node {}}
          };
        
        \node[draw=none, right] at (2.5,0) {\small $k=0 \rightarrow 2^0 = 1$ nó};
        \node[draw=none, right] at (2.5,-0.8) {\small $k=1 \rightarrow 2^1 = 2$ nós};
        \node[draw=none, right] at (2.5,-1.6) {\small $k=2 \rightarrow 2^2 = 4$ nós};
        
        \draw[dashed, gray] (root) -- (2.5,0);
        \draw[dashed, gray] (1.5,-0.8) -- (2.5,-0.8);
        \draw[dashed, gray] (2.25,-1.6) -- (2.5,-1.6);
    \end{tikzpicture}
\end{center}
\begin{enumerate}[label=\alph*)]
    \item Utilizando a fórmula apropriada, quantos nós \textit{no máximo} podem existir no nível $k=66$ (considerando que a raiz ocupa o nível 0)?
    \item Se uma árvore binária é perfeita e possui altura $h=312$, qual o seu número total de nós? Demonstre o cálculo.
\end{enumerate}
\quadro{7cm}

% -------------------------------------------------------
% Questão 10: Lema e Pior Caso
% -------------------------------------------------------
\item \textbf{Pior Caso (Degeneração) e a Relação das Arestas}
\begin{enumerate}[label=\alph*)]
    \item Desenhe uma árvore degenerada com 5 nós. Por que cientistas da computação consideram essa formação um ``pior caso'' para a busca (qual a sua notação Big-O)?
    \item Baseando-se no Lema do Aperto de Mãos aplicado a árvores (Grafos Acíclicos Conexos), quantas arestas obrigatoriamente existem em uma árvore binária que possui exatos 150 vértices? Justifique matematicamente porque não pode haver uma aresta a mais nem uma a menos.
\end{enumerate}
\quadro{10cm}

\end{enumerate}

\end{document}
